\newpage
\section{Teoria di Hamilton--Jacobi}
\note{Lezione 05/05/2026}

Ci poniamo il seguente problema:
\begin{quote}
    Determinare una trasformazione canonica $Q_i=Q_i(\vq,\vp,t)$, $P_i=P_i(\vq,\vp,t)$ tale che la nuova hamiltoniana $K$ sia identicamente nulla, cioé $K \equiv 0$.
\end{quote}
Con una trasformazione canonica che soddisfa questa condizione, le equazioni di Hamilton diventano 
\begin{equation*}
    \begin{cases}
        \dot{Q}_i = \pdv{K}{P_i} = 0 \\
        \dot{P}_i = -\pdv{K}{Q_i} = 0
    \end{cases}
\end{equation*}
e dunque $Q_i$ e $P_i$ sono costanti di moto
\begin{equation*}
    Q_i = \text{costante} = \beta_i, \quad P_i = \text{costante} = \alpha_i.
\end{equation*}

Determiniamo una tale trasformazione. Sappiamo che una trasformazione canonica può essere ottenuta mediante una funzione generatrice $F_2(\vq,\vP,t)$, con le seguenti relazioni:
\begin{equation*}
    \begin{cases}
        p_i = \pdv{F_2}{q_i} \\
        Q_i = \pdv{F_2}{P_i} \\
        K = H + \pdv{F_2}{t}
    \end{cases}
\end{equation*}
posto $\det \left( \pdv{F_2}{q_i}{P_j} \right) \neq 0$ per garantire l'invertibilità della trasformazione. La condizione $K \equiv 0$ implica che
\begin{equation}
    \label{eq:hamiltonjacobi-1}
    H + \pdv{F_2}{t} = 0.
\end{equation}
L'hamiltoniana $H$ è una funzione di $\vq$, $\vp$ e $t$, ma possiamo esprimere $\vp$ in funzione di $\vq$, $\vP$ e $t$ tramite la prima relazione della trasformazione, ottenendo così un'equazione differenziale parziale per $F_2(\vq,\vP,t)$:
\begin{equation*}
    H\left( \vq, \pdv{F_2}{\vq}, t \right) + \pdv{F_2}{t} = 0.
\end{equation*}

In letteratura la funzione $F_2$ viene indicata con $S$ e viene detta \emph{funzione principale di Hamilton}
\begin{align*}
    S(\vq,\vP,t) &\defeq F_2(\vq,\vP,t).
\end{align*}
Poiché $P_i = \alpha_i$ rappresenta una costante di moto, possiamo riscrivere $S$ come funzione delle sole costanti arbitrarie piuttosto che dei momenti coniugati:
\begin{equation*}
    S(\vq,\vP,t) = S(\vq,\alpha_1, \dots, \alpha_n, t).
\end{equation*}
La trasformazione canonica associata assume quindi la forma
\begin{equation*}
    \begin{cases}
        p_i = \pdv{S}{q_i} \\
        \beta_i = \pdv{S}{\alpha_i}
    \end{cases}
\end{equation*}
Sostituendo questa espressione nella condizione \eqref{eq:hamiltonjacobi-1}, ricaviamo
\begin{equation*}
    H\left( q_1, \dots, q_n, \pdv{S}{q_1}, \dots, \pdv{S}{q_n}, t \right) + \pdv{S}{t} = 0,
\end{equation*}
che prende il nome di \emph{equazione di Hamilton--Jacobi}.

\subsection{Soluzioni dell'equazione di Hamilton--Jacobi}

A differenza delle equazioni di Hamilton, che sono equazioni differenziali ordinarie, l'equazione di Hamilton--Jacobi è un'equazione differenziale alle derivate parziali di primo ordine nella funzione incognita $S$. Il problema diventa pertanto determinare soluzioni di questa EDP, ricordando che la soluzione generale dipende genericamente da funzioni arbitrarie.

Consideriamo l'equazione di Hamilton--Jacobi
\begin{equation*}
    H(q_1, \dots, q_n, \pdv{S}{q_1}, \dots, \pdv{S}{q_n}, t) + \pdv{S}{t} = 0.
\end{equation*}
Calcoliamo la derivata totale di $S$ rispetto al tempo
\begin{equation*}
    \dv{S}{t} = \sum_{i=1}^n \pdv{S}{q_i} \dot{q}_i + \pdv{S}{t} = \sum_{i=1}^n p_i \dot{q}_i - H = L,
\end{equation*}
ritrovando così la lagrangiana $L$. Potrebbe sembrare che il problema sia facilmente risolvible integrando
\begin{equation*}
    S = \int L(\vq, \vp, t) \dd{t} + \text{costante},\note{La funzione principale di Hamilton è proprio l'azione}
\end{equation*}
tuttavia $L$ è una funzione stessa delle soluzioni, tornando al punto di partenza.

Cerchiamo allora di trovare una sottoclasse di soluzioni.

\begin{definition}[Integrale completo dell'equazione di Hamilton--Jacobi]
    Una funzione $S(\vq,\alpha_1, \dots, \alpha_n, t)$, dove $\alpha_1, \dots, \alpha_n$ sono costanti arbitrarie, si dice \emph{integrale completo} o \emph{soluzione completa} dell'equazione di Hamilton--Jacobi se è soluzione di tale equazione per ogni scelta delle $\alpha_i$ e verifica la condizione
    \begin{equation*}
        \det \left( \pdv{S}{q_i}{\alpha_j} \right) \neq 0.
    \end{equation*}
\end{definition}

\begin{theorem}[Jacobi]
    Se $S(\vq,\alpha_1, \dots, \alpha_n, t)$ è una soluzione completa dell'equazione di Hamilton--Jacobi, allora la soluzione generale delle equazioni di Hamilton è data dalle seguenti:
    \begin{equation*}
        p_i = \pdv{S}{q_i}, \quad \beta_i = \pdv{S}{\alpha_i}.
    \end{equation*}
\end{theorem}

\begin{remark}
    Una soluzione completa dell'equazione di Hamilton--Jacobi non è la soluzione generale di tale equazione, bensì una sottoclasse particolare caratterizzata dalla presenza di $n$ costanti arbitrarie.
\end{remark}


Per inquadrare correttamente il ruolo delle costanti arbitrarie in questa teoria, consideriamo una generica equazione differenziale alle derivate parziali di primo ordine in $n+1$ variabili indipendenti
\begin{equation*}
    F(x_1, \dots, x_{n+1}, \pdv{f}{x_1}, \dots, \pdv{f}{x_{n+1}}) = 0,
\end{equation*}
la soluzione completa dipende da $n+1$ costanti arbitrarie. Nel nostro caso, tuttavia, il tempo non figura come variabile indipendente alla stessa stregua delle coordinate spaziali nell'hamiltoniana. Pertanto, la $(n+1)$-esima costante si manifesta come una costante additiva in $S$, che risulta fisicamente ininfluente poiché le trasformazioni canoniche dipendono da $\pdv{S}{q_i}$ e $\pdv{S}{\alpha_i}$, annullando tale contributo.

Nel caso di un sistema con hamiltoniana indipendente dal tempo si ha
\begin{equation*}
    \pdv{H}{t} = 0 \implies H(\vq, \vp, t) = \text{costante} = \gamma,\note{Si noti che $\gamma$ è una costante solo rispetto a tempo}
\end{equation*}
e integrando $\dv{S}{t} = -\gamma$, si può cercare una soluzione del tipo
\begin{equation*}
    S(\vq,\alpha_1, \dots, \alpha_n, t) = W(\vq,\alpha_1, \dots, \alpha_n) - \gamma t,
\end{equation*}
dove in generale $\gamma$ dipende dalle $\alpha_i$ e $W$ è detta \emph{funzione caratteristica di Hamilton}. Sostituendo questa forma nell'equazione di Hamilton--Jacobi si ottiene
\begin{equation*}
    H\left( \vq, \pdv{W}{q_1}, \dots, \pdv{W}{q_n} \right) = \gamma,
\end{equation*}
la quale è un'equazione differenziale alle derivate parziali di primo ordine per $W$, detta \emph{equazione ridotta di Hamilton--Jacobi}.

In questo contesto, è naturale scegliere la costante $\gamma$ come una delle costanti arbitrarie $\alpha_i$; ad esempio, ponendo $\gamma = \alpha_1$. La ridotta equazione di Hamilton--Jacobi assume allora la forma
\begin{equation*}
    H\left( \vq, \pdv{W}{q_1}, \dots, \pdv{W}{q_n} \right) = \alpha_1,
\end{equation*}
che risulta particolarmente utile nelle applicazioni pratiche.

\begin{example}[Applicazione della teoria di Hamilton--Jacobi all'oscillatore armonico]
    Nel caso dell'oscillatore armonico unidimensionale, l'hamiltoniana è indipendente dal tempo ed è proprio uguale all'energia meccanica $\alpha$
    \begin{equation*}
        H(q,p) = \frac{1}{2m}\left( p^2 + m^2 \omega^2 q^2 \right)= \alpha,
    \end{equation*}
    dove $\alpha$ è una costante di moto. Pertanto si ha
    \begin{equation*}
        S = W(q) - \alpha t,
    \end{equation*}
    con le trasformazioni canoniche
    \begin{equation}
        \label{eq:oscillatoreHamiltonJacobi}
        p = \pdv{W}{q}, \quad \beta' = \pdv{W}{\alpha} - t.
    \end{equation}

    L'equazione ridotta di Hamilton--Jacobi diventa
    \begin{equation*}
        \frac{1}{2m}\left[ \left( \pdv{W}{q} \right)^2 + m^2 \omega^2 q^2 \right] = \alpha,
    \end{equation*}
    da cui si ricava
    \begin{equation*}
        \pdv{W}{q} = \sqrt{2m\alpha - m^2 \omega^2 q^2}
    \end{equation*}
    e integrando si ottiene la soluzione
    \begin{align*}
        W(q,\alpha) &= \int \sqrt{2m\alpha - m^2 \omega^2 q^2} \dd{q}.
    \end{align*}
    Deriviamo rispetto ad $\alpha$
    \begin{align*}
        \pdv{W}{\alpha} &= \int \frac{m}{\sqrt{2m\alpha - m^2 \omega^2 q^2}} \dd{q} \\
        &= \frac{m}{\sqrt{2m\alpha}} \int \frac{1}{\sqrt{1 - \frac{m \omega^2 q^2}{2\alpha}}} \dd{q} \\
        &= \sqrt{\frac{m}{2\alpha}} \sqrt{\frac{2\alpha}{m\omega^2}} \arcsin \left( q\sqrt{\frac{m\omega^2}{2\alpha}} \right) \\
        &= \frac{1}{\omega} \arcsin \left( q\sqrt{\frac{m\omega^2}{2\alpha}} \right).
    \end{align*}
    Sostituendo questa espressione nella seconda relazione delle trasformazioni canoniche \eqref{eq:oscillatoreHamiltonJacobi} si ottiene
    \begin{equation*}
        \arcsin \left( q\sqrt{\frac{m\omega^2}{2\alpha}} \right) = \omega t + \omega \beta'.
    \end{equation*}
    Ponendo $\omega \beta' = \beta$ e applicando la funzione seno ad entrambi i membri si ricava la soluzione generale dell'oscillatore armonico
    \begin{equation*}
        q(t) = \sqrt{\frac{2\alpha}{m\omega^2}} \sin(\omega t + \beta).
    \end{equation*}

    Calcolando $p$ tramite la prima relazione delle trasformazioni canoniche \eqref{eq:oscillatoreHamiltonJacobi} si ottiene
    \begin{align*}
        p(t) &= \sqrt{2m\alpha} \sqrt{1- \frac{m\omega^2}{2\alpha} q^2} \\
        &= \sqrt{2m\alpha} \sqrt{1- \sin^2(\omega t + \beta)} \\
        &= \sqrt{2m\alpha} \cos(\omega t + \beta).
    \end{align*}
    
    Otteniamo facilmente le condizioni iniziali
    \begin{equation*}
        q_0 = \sqrt{\frac{2\alpha}{m\omega^2}} \sin(\beta), \quad p_0 = \sqrt{2m\alpha} \cos(\beta)
    \end{equation*}
    e quindi, invertendo le relazioni, si ricavano le costanti arbitrarie $\alpha$ e $\beta$ in funzione delle condizioni iniziali
    \begin{equation*}
        \alpha = \frac{p_0^2}{2m} + \frac{1}{2} m \omega^2 q_0^2, \quad \beta = \arctan \left( \frac{m\omega q_0}{p_0} \right).
    \end{equation*}
\end{example}